\section{总结与延伸}

\subsection{核心收获}

LangChain 这篇文章传递了一个强有力的信息：\textbf{在当前 AI Agent 的实践中，Harness Engineering 的边际收益可能远高于等待更强模型}。

关键数据支撑：
\begin{itemize}[leftmargin=1.5em]
    \item 同一模型（GPT-5.2-Codex），仅改 Harness：52.8\% $\rightarrow$ 66.5\%（+13.7pp）
    \item 排名从 Top 30 跃升至 Top 5
    \item 改进手段全部是"系统工程"层面的——Prompt、Middleware、上下文注入
\end{itemize}

\subsection{对 Agent 从业者的启示}

\begin{importantbox}{Harness Engineering 的投资回报}
在构建 Agent 应用时，不要只关注"用哪个模型"。更应该投入时间在：
\begin{itemize}[leftmargin=1.5em]
    \item 设计 Trace 驱动的迭代改进流程
    \item 构建 Self-Verification 机制
    \item 做好 Context Engineering
    \item 识别并修补模型的行为短板
\end{itemize}
这些投入的性价比往往远高于简单地换一个更大/更贵的模型。
\end{importantbox}

\subsection{未来展望}

LangChain 指出了几个值得关注的前沿方向：
\begin{enumerate}[leftmargin=2em]
    \item \textbf{多模型系统}：让 Codex、Gemini、Claude 协同工作，各取所长
    \item \textbf{持续学习的 Memory 原语}：让 Agent 能跨任务积累经验并自主改进
    \item \textbf{跨模型 Harness 评估}：系统性地度量 Harness 变更对不同模型的影响
    \item \textbf{RLM（Reinforcement Learning from Mistakes）}：更高效地从 Trace 中挖掘改进信号
\end{enumerate}

此外，LangChain 已经开源了 Deep Agents（提供 Python 和 JavaScript 版本），并公开了实验 Trace 数据集供社区研究。

\subsection{拓展阅读}

\begin{itemize}[leftmargin=1.5em]
    \item LangChain Blog 原文：\href{https://blog.langchain.com/improving-deep-agents-with-harness-engineering/}{Improving Deep Agents with Harness Engineering}
    \item Deep Agents 开源仓库（Python \& JavaScript）
    \item LangSmith：LangChain 的 Trace 和评估平台
    \item Terminal Bench 2.0：Agentic Coding 基准测试
    \item Harbor：Agent 评估编排工具
\end{itemize}

%% --- Fin del contenido --- %%

\end{document}
